% CAL-13: materiał kanoniczny w ME-20, ME-21, ME-32 i FA-12.
% Wszystkie 44 etykiety oryginalnego rozdziału zachowano w tych źródłach.
% Do wycofania książki pozostawiamy jedynie odsyłacze.
\providecommand{\atlasref}[2]{#2}
\section{Całka Radona}

Topologiczne własności miar i reprezentację dodatnich
funkcjonałów przez całki omawiamy w teorii miary.
Nie powtarzamy tutaj definicji, dowodów ani przykładów.

\begin{itemize}
 \item \atlasref{ME-20::def:radon-measure-20}{Miary Radona i regularność};
 \item \atlasref{ME-20::lem:cutoff-radon-20}{Funkcje odcinające o zwartym nośniku};
 \item \atlasref{ME-20::thm:uniqueness-cc-20}{Jednoznaczność miary Radona};
 \item \atlasref{ME-21::thm:riesz-positive-21}{Twierdzenie Riesza--Markowa--Kakutaniego dla lokalnie zwartych przestrzeni Hausdorffa};
 \item \atlasref{ME-21::thm:compact-riesz-21}{Alternatywna konstrukcja na zwartych przestrzeniach metrycznych};
 \item \atlasref{ME-21::eq:gluing-measures-21}{Sklejanie miar na wyczerpaniu};
 \item \atlasref{ME-21::thm:riesz-c0-positive-21}{Dodatnie funkcjonały na przestrzeni $C_0(X)$};
 \item \atlasref{ME-21::thm:riesz-complex-21}{Reprezentacja funkcjonałów zespolonych};
 \item \atlasref{FA-12::prop:ck-banach12}{Zupełność przestrzeni $C(K)$};
 \item \atlasref{ME-32::it:thm:daniell-radon-comparison}{Porównanie miar Daniella i Radona};
 \item \atlasref{ME-32::it:cor:radon-cc-l1-general}{Gęstość funkcji o zwartym nośniku w $L^1$}.
\end{itemize}
Przykłady miar atomowych, miar skupionych na okręgu
oraz sum miar atomowych i miary Lebesgue'a znajdują się w
\atlasref{ME-21::ch:riesz-markov-21}{rozdziale 21 teorii miary}.
